\documentclass[11pt,letterpaper]{article}
\usepackage[T1]{fontenc}
\usepackage{lmodern,amsmath,amssymb,amsthm,booktabs,array,geometry,xcolor,fancyhdr,microtype}
\usepackage{hyperref}
\geometry{left=0.92in,right=0.92in,top=0.88in,bottom=0.9in}
\definecolor{ink}{HTML}{194955}
\definecolor{muted}{HTML}{48636A}
\hypersetup{colorlinks=true,urlcolor=ink,linkcolor=ink,citecolor=ink,
 pdftitle={A counterexample to McNew's conjecture on geometric progressions},
 pdfauthor={Research report}}
\pagestyle{fancy}\fancyhf{}\fancyhead[L]{\footnotesize\sffamily\color{muted}METHODS FROM THE LEAN PROOF OF DIRICHLET GRH}
\fancyhead[R]{\footnotesize\sffamily\color{muted}26 NOVEMBER 2026}
\fancyfoot[C]{\small\thepage}\setlength{\headheight}{14pt}
\renewcommand{\headrulewidth}{0.25pt}
\newtheorem{theorem}{Theorem}[section]
\newtheorem{corollary}[theorem]{Corollary}
\newcommand{\upsrc}[2]{\href{https://github.com/lsernau/grh-lean/blob/72fead102f663d5c7a407964a13dc0cf9d977207/grh/#1}{#2}}
\newcommand{\prsrc}[2]{\href{https://github.com/VladimirReshetnikov/ProveIt/blob/a1d43d8a32add57d596273e04648bc96a0a9bc0d/docs/incoming/#1}{#2}}
\setlength{\parindent}{1.2em}
\allowdisplaybreaks
\begin{document}
\thispagestyle{empty}
\noindent\textcolor{ink}{\small\sffamily RESEARCH REPORT}\hfill {\small\sffamily 26 November 2026}
\par\vspace{8pt}\noindent{\color{ink}\rule{\linewidth}{0.8pt}}\par\vspace{7pt}
\noindent{\fontsize{26}{29}\selectfont\bfseries A counterexample to McNew's\\geometric-progression conjecture}\par\vspace{4pt}
\noindent{\large Further applications of methods in the Lean proof of Dirichlet GRH}\par\vspace{9pt}

\noindent Rankin constructed a set of positive integers of natural density about $0.71974$ that avoids every nonconstant three-term geometric progression of rational ratio. McNew conjectured that this density is the largest possible \cite[Conj.~7.2]{mcnew}. We construct a set of strictly larger natural density. The proof marks all added integers with the same forbidden exponent at $2$ and makes a different odd prime carry the other forbidden exponent. A single family of deletions protects all these additions. Their reciprocal mass grows like a sum of reciprocals of primes, while the extra deletions cost a convergent sum of inverse prime squares. A second construction on primes through $1427$ gives a natural-density gain exceeding $10^{-596}$. The same construction improves the natural norm-density of the analogous set of ideals in every number field.

\paragraph{A separate upper-density bound.} For the supremum $\bar\alpha$ of \emph{upper} densities, a construction with 275 smooth-integer templates and seven finite certificates gives
\[
                  0.7365954<\bar\alpha<0.7716011.
\]
Both endpoints improve on the interval in McNew's paper. The two assertions concern different extrema: $d(A)$ in the counterexample is a limit, while the template construction controls a limsup.

\paragraph{An exact answer for cyclotomic witnesses.}
A prior ProveIt manuscript asks for the greatest common divisor of \emph{all} positive witnesses in an intermediate annihilation theorem. We determine it for every modulus $N$ and weight $e\ge1$. The character contribution is the common denominator of $1/(2L(-e,\chi))$ in the character rings. Passing to a larger modulus before canceling the factorial in Voronoi's congruence is the step that makes the denominator sufficient. At $(N,e)=(3,9)$ this excludes the coefficient primes $5$ and $7$ from the support of the specified odd components.

\paragraph{Two structural results.}
We determine the largest set of sparse places allowed by the two prime constraints at $2$ and $3$ in the source. For any fixed prime cutoff at least three, we obtain the sharp asymptotic order of its size. On the analytic side, we classify binary Frobenius tests for which the undesired factors have zero coefficient and all remaining nonabelian coefficients have one prescribed sign (allowing arbitrary linear coefficients). For $X^3-X-1$ this yields a density-$1/4$ event testing exactly one quadratic twist of the cubic's standard representation. RH for that twist is equivalent to a precise upper estimate for the event, conditional only on the upstream Dirichlet GRH theorem. That upper estimate remains open.

\paragraph{Status of the proofs.} The density bounds, the cyclotomic formula, the sparse-place estimates and the finite-group classification are independent of RH. Sections~\ref{sec:rankin}--\ref{sec:binary} contain proofs, and the archive contains exact finite verifiers. These extensions are written mathematics. The Lean verification reported in \cite{grh} pertains to the upstream theorems. Priority claims here are limited to explicit comparisons with the cited papers and manuscripts.
\clearpage
\section{A counterexample to McNew's natural-density conjecture}\label{sec:rankin}
For $E\subseteq\mathbb N$, its natural density is $d(E)=\lim_{x\to\infty}|E\cap[1,x]|/x$ when the limit exists. Let $\alpha$ be the supremum of these densities over sets avoiding nonconstant rational-ratio three-term progressions. McNew conjectured that Rankin's set attains $\alpha$ \cite[Conj.~7.2]{mcnew}. Let $G$ be the nonnegative integers whose ternary digits are all $0$ or
$1$. Rankin's set is $R=\{n\ge1:v_p(n)\in G\text{ for every prime }p\}$.
The set $G$ has no nonconstant three-term arithmetic progression:
in $u+w=2v$, the lowest ternary digit at which the three values differ
would give a nonconstant solution with digits in $\{0,1\}$ modulo $3$.
The valuations of a rational-ratio geometric progression satisfy this
identity at every prime. Thus $R$ is progression-free.
It has natural density approximately $0.71974$. Riddell had improved the corresponding constructions for lengths at least four \cite{riddell}. Here the length is three.

\medskip\noindent\textbf{Theorem.}
There is a subset $A$ of the positive integers with a natural density
strictly greater than $d(R)$ and with no nonconstant three-term
geometric progression of rational ratio. In particular
$\alpha>d(R)$, which disproves Conjecture 7.2 of McNew's published paper.

\medskip\noindent\textit{Construction.}
Choose a finite set $B$ of odd primes such that
$\sum_{q\in B}1/q>13$. One definite choice is the shortest initial
segment of the odd primes with this property. Euler's divergence theorem
ensures that the segment is finite. Alternatively use every odd prime up
to the explicit integer $Y_0=\lceil\exp(\exp(15))\rceil$. The finite Euler
product over $p\le Y_0$ is at least $\sum_{n\le Y_0}1/n\ge\log Y_0$. In its
logarithm the prime powers of exponent at least two contribute at most
$\sum_{n\ge2}1/[n(n-1)]=1$. Thus the odd-prime reciprocal sum up to
$Y_0$ is at least $\log\log Y_0-1-1/2\ge13.5$. Put $P=2\prod_{q\in B}q$ and
$S=\{2\}\cup B$. For the moment work solely with the semigroup of
$S$-smooth positive integers. Write $C$ for Rankin's smooth integers.
Make the following deletions from $C$.
\begin{itemize}
\item Delete those whose exponent equals $1$ at every prime of $B$.
\item For each $q\in B$, delete those whose exponent is at least $3$ at
$q$ and equals $1$ at every other prime of $B$.
\end{itemize}
In both rules the exponent at $2$ can be any member of $G$.
To the survivors add $X=\{2Pq:q\in B\}$ and call the resulting smooth
set $F$. Each added integer has exponent $2$ at $2$ and at exactly one
odd prime $q$, and exponent $1$ at all other odd primes in $B$.

\medskip\noindent\textit{Progression avoidance.}
Three retained old points were already in $C$. Suppose a progression
contained exactly one new point, $2Pq$. At every $r\in B\setminus\{q\}$
the three exponents belong to $G$ and one equals $1$, so all three equal
$1$. In this slice the deletion rules leave an old point with exponent
at $q$ only if that exponent is $0$: exponent $1$ was deleted by the
first rule, $2\notin G$, and exponents at least $3$ were deleted by the
second. The three $q$-exponents would therefore be $0,0,2$ in some
order, which cannot be an arithmetic progression. With two new points,
two $2$-adic exponents equal $2$, forcing the third to equal $2$;
an old point cannot do so. Three new points can form no nonconstant
progression because every odd coordinate lies in $\{1,2\}$ and the
$2$-adic coordinate is constant. These cases permit any ordering of the
three terms and therefore cover all rational ratios.

\medskip\noindent\textit{Gain in reciprocal mass.}
For a set of $S$-smooth integers put $W(E)=\sum_{n\in E}1/n$;
the sum converges because $S$ is finite.
The added mass is $P^{-1}\sum_{q\in B}1/(2q)$.
For deleted points of the first type the odd part is $P/2$,
so their mass in units of $1/P$ is
$2\sum_{e\in G}2^{-e}\le4$. For the second type, raising the
$q$-exponent from $1$ to $e\ge3$ multiplies the reciprocal weight
by $q^{1-e}$. Summing over all $e\ge3$ gives $1/[q(q-1)]$.
It follows that
\[
 P W(C\setminus F)
 \le 4\left(1+\sum_{q\in B}\frac{1}{q(q-1)}\right)
 \le 4\left(1+\sum_{n=3}^\infty\frac{1}{n(n-1)}\right)=6.
\]
No new point belongs to $C$. Consequently
\[
                 W(F)-W(C)>\frac{1}{2P}.
\]
The saving comes from paying once for the deleted slice with all odd
exponents equal to $1$. Each new prime contributes $1/(2Pq)$, while
the additional deletion cost at that prime is bounded by
$4/[Pq(q-1)]$. Requiring every addition to have $2$-adic exponent
exactly $2$ controls progressions involving two additions simultaneously.

\medskip\noindent\textit{Natural density.}
Let $H$ be the positive integers coprime to $P$ with $v_p(h)\in G$
at every prime outside $S$. They have positive natural density
\[
 c_H=\prod_{p\in S}(1-p^{-1})
      \prod_{p\notin S}\left((1-p^{-1})\sum_{e\in G}p^{-e}\right),
\]
where the second product multiplies the whole local factor
$(1-p^{-1})\sum_{e\in G}p^{-e}$ at each $p\notin S$.
Each factor in the second product is at least $1-p^{-2}$, since $0,1\in G$.
Finite-prime valuation frequencies follow from the Chinese remainder
theorem by truncating the exponents. Omitting conditions at primes
$p>Y$ costs upper density at most $\sum_{p>Y}p^{-2}\to0$, proving the
product formula. Set $A=\{hf:h\in H,\ f\in F\}$. In any progression
the valuations outside $S$ are constant by the progression-freeness of
$G$, so the three rough parts $h$ agree. Cancelling $h$ reduces the
progression to one in $F$. This proves the global avoidance claim.

By unique factorization the sets $fH$ are disjoint for distinct smooth
$f$, and $d(fH)=c_H/f$. For each $K$ the union over $f>K$ contributes
upper density at most $\sum_{f>K}1/f$, tending to zero. Hence $A$ has
natural density $c_H W(F)$, while $d(R)=c_HW(C)$ by the identical
argument. In particular
\[
                         d(A)-d(R)>\frac{c_H}{2P}>0 .
\]
For the explicit cutoff just given, $c_H\ge(\varphi(P)/P)/\zeta(2)>1/(2P)$
and $P\le Y_0!$, so the gain even exceeds $1/[4(Y_0!)^2]$.
The bound proves strict improvement; its deliberately large finite
prime set makes it unsuitable for a visible change in the decimal
approximation $0.71974$. The greatest possible natural density remains
to be determined. This proof uses Euler's theorem and elementary
valuation arguments; no computational search or zero theorem enters it.

\paragraph{Why optimizing a fixed finite list of primes can miss the improvement.}
Fix a finite set of primes $T$, a nonempty progression-free family $D$ of $T$-smooth integers, and Rankin's original rule at all primes outside $T$. The resulting product construction always admits a strict improvement. Indeed, choose a marker prime $\ell\notin T$ and put all satellites $q\in B$ outside $T\cup\{\ell\}$. Replacing $2$ by $\ell$ in the proof makes the normalized added mass $\ell^{-1}\sum_{q\in B}q^{-1}$ and bounds the deleted mass by $\ell(1-1/\ell)^{-1}(1+\sum_{q\in B}1/[q(q-1)])$. The first quantity is unbounded as $B$ grows. The two unchanged factors, namely $D$ and the remaining unmodified primes, multiply the gain by a positive constant in the density formula. A progression would project to a progression in each block, forcing constancy in every block. All changed integers can also be required to exceed any prescribed cutoff by choosing $B$ beyond it. This explains how exact finite-prime optimality results can coexist with failure of the unrestricted conjecture.

\paragraph{Relation to the previous approaches.}
Rankin's independent prime-exponent rule, McNew's comparisons between ratio conventions, and the prior report's sharp four-prime result all use the same valuation identity. The Lean proof's sparse-place argument is another instance of converting multiplicative patterns to forbidden additive relations. The step here is to \emph{correlate} prime exponents: fix one forbidden marker and share a single deletion among additions indexed by many primes. No estimate for zeros or theta lifts is required.

\subsection{A finite counterexample with a certified gain}

The counterexample can use just the $224$ odd primes $q\le Y=1427$.
Let $B$ denote these primes, $Q=\prod_{q\in B}q$, $S=\{2\}\cup B$, and
$C$ the Rankin subset of the $S$-smooth integers. Write $a=v_2(n)$ and
$b_q=v_q(n)$. For an old point $n\in C$, delete it under either of the
following conditions.
\begin{itemize}
\item All its odd exponents $b_q$ equal $1$ and $a>0$.
\item Exactly one odd exponent $b_q$ differs from $1$, and either $b_q=4$,
or $b_q=3$ and $a\in2K$, where $K=2+3\{0,1,3,4\}=\{2,5,11,14\}$ and $2K=\{2k:k\in K\}$.
\end{itemize}
Insert the points $2^kQ$, $2^kQq$ and $2^kQq^3$ for every $q\in B$ and $k\in K$.
All four markers are $2\bmod 3$, so these points are disjoint from $C$.

\paragraph{The local obstruction.}
Let $G$ be the old ternary exponent set.
If $u,v\in G$ and $u,2,v$ is an arithmetic progression, then
$\{u,v\}=\{0,4\}$ or $\{1,3\}$. If $u,v,2$ is a progression, then
$u=2v-2$ and $v\in\{1,3\}$. To verify the last assertion, write $v$
in ternary. If the last digit is $1$, subtracting $2$ from $2v$ leaves
each higher digit doubled, so all higher digits must vanish. If the
last digit is $0$, let $j\ge1$ be the position of the first nonzero
digit. Borrowing to subtract $2$ would create a digit $2$ at position
$1$ if $j>1$. Thus $j=1$ and all higher digits must vanish, again by
doubling. The two possible $v$ do work.

Consider a progression with only the insertion $2^kQq$ among its three
members. All its other odd exponents are $1$, by progression avoidance
in $G$. If the insertion is an endpoint, the middle $q$-exponent is
$1$ or $3$ by the local obstruction. In the first case, the middle
point has all odd exponents $1$ and positive $2$-exponent, since it
averages the positive $k$ and a nonnegative exponent. In the second
case, the other endpoint has $q$-exponent $4$. Both cases give a
deleted point. If the insertion is the midpoint, the outside
$q$-exponents are $0,4$ or $1,3$. The point with $q$-exponent $4$
was deleted. In the other pair the point with $q$-exponent $1$ survives
only at $a=0$, which forces the point with $q$-exponent $3$ to have
$a=2k\in2K$, also deleted. For a single insertion $2^kQ$, both
old points must have all odd exponents $1$, and can survive only at
$a=0$. The values $0,0,k$ cannot form a progression. For a single insertion $2^kQq^3$, the exceptional odd exponent is $4\in G$. Every odd coordinate of all three points then belongs to $G$ and is constant. Both old points would have exactly the deleted odd pattern with exceptional exponent $4$.

If two points were inserted, their $2$-exponents are $2\bmod3$, forcing
the third exponent to have that residue in any of the three positions.
No member of $G$ has this residue. With three insertions, every odd exponent
is in the progression-free set $\{1,2,4\}$ and the $2$-exponent is in the progression-free set $K$ (an affine image of a subset of $G$). Each coordinate
must therefore be constant in a progression, so all terms coincide.
This proves avoidance for every rational ratio.

\paragraph{Size of the improvement.}
Let $S_2=\sum_{a\in G}2^{-a}$. Exponents below $27$ use only the
three lowest ternary digits. Bounding the remaining terms by
$\sum_{a\ge27}2^{-a}=2^{-26}$ gives
\[
 S_2\le U=(1+2^{-1})(1+2^{-3})(1+2^{-9})+2^{-26}.
\]
Write $u=\sum_{k\in K}2^{-k}$ and $v=\sum_{k\in K}2^{-2k}$.
Let $\Delta W$ denote the change in reciprocal mass on smooth numbers.
In units of $1/Q$, the insertions contribute $u(1+\sum_{q\in B}(q^{-1}+q^{-3}))$,
the common deletion costs $S_2-1$, and the two kinds of
single-coordinate deletions cost at most $v/q^2$ and exactly $S_2/q^3$,
respectively. The first estimate also counts any $2k$ absent from $G$. Consequently
\[
 Q\Delta W\ \ge\ u-(U-1)
     +\sum_{q\in B}\left(\frac{u}{q}-\frac{v}{q^2}-\frac{U-u}{q^3}\right)
       >\frac9{100000}.
\]
The last strict inequality is checked in exact rational arithmetic by
\texttt{check\_rankin\_finite.py}. The script sieves the primes, verifies
$Q<3\cdot10^{590}$, and checks
\[
          \frac{Y}{Y+1}\prod_{p\le Y}(1-p^{-1})>\frac1{15}.
\]
The outside-prime density factors are at least $1-p^{-2}$, so
their product is at least $\prod_{n>Y}(1-n^{-2})=Y/(Y+1)$.
The density $c_H$ of the permitted rough parts is therefore greater
than $1/15$. Multiplying the smooth gain by $c_H$ yields the explicit
bound for the natural density of the full construction,
\[
              d(A)-d(R) > \frac{9}{100000\cdot15\cdot3\cdot10^{590}}
                           > 10^{-596}.
\]

\subsection*{A consequence for integral ideals}
In every number field $K$, the set of nonzero integral ideals whose
prime-ideal exponents belong to $G$ admits a strict improvement in
natural density by norm, preserving avoidance even for fractional-ideal
ratios. This gives a qualitative answer to the natural-density aspect
of Best, Huan, McNew, Miller, Powell, Tor and Weinstein's Question 5.1
for quadratic fields
(\href{https://arxiv.org/abs/1412.0999}{their paper, section 5}).
Their question asks how much nongreedy sets can improve existing bounds.
This argument gives a positive improvement on the greedy lower bound
for natural density. Sharp field-specific improvements and the extremal
upper density are separate questions.

Indeed, fix a prime ideal $\mathfrak l$ of norm $Q$ and choose a finite
set $B$ of other prime ideals. In the preceding construction replace
$2$ by $\mathfrak l$, $q$ by $\mathfrak q$, and $P$ by
$\mathfrak l\prod_{\mathfrak q\in B}\mathfrak q$. Every combinatorial
step is unchanged by unique factorization of ideals. Give an ideal
reciprocal mass $1/N\mathfrak a$. In units of $1/NP$, the additions
have weight $Q^{-1}\sum_{\mathfrak q\in B}1/N\mathfrak q$, while the
deletions have weight at most
\[
  \frac{Q}{1-Q^{-1}}
  \left(1+\sum_{\mathfrak q\in B}
    \frac1{N\mathfrak q\,(N\mathfrak q-1)}\right).
\]
The sum in parentheses is bounded uniformly in $B$ since
$\sum_{\mathfrak q}(N\mathfrak q)^{-2}<\infty$.
In contrast, $\sum_{\mathfrak q}1/N\mathfrak q$ diverges:
the logarithm of the Euler product for $\zeta_K(s)$ tends to infinity
as real $s\downarrow1$, and terms from prime-ideal powers of exponent
at least two have bounded total. Thus the gain is positive for some
finite $B$.

To justify natural norm-density, the count $J_K(x)$ of integral ideals
of norm at most $x$ satisfies $J_K(x)\sim c_Kx$ with $c_K>0$
(\href{https://doi.org/10.1016/j.exmath.2006.08.001}{Murty and Van Order}).
Consequently multiples of an ideal $\mathfrak d$ have density
$1/N\mathfrak d$, and $J_K(x)\ll_K x$ for all $x\ge1$.
Inclusion--exclusion yields the local factors for finitely many prime
ideals. The $\mathfrak q^2$-multiple tail has upper density bounded
by a constant times $\sum_{N\mathfrak q>Y}(N\mathfrak q)^{-2}$.
The global density, its positivity and the strict improvement now
follow by the same convergent smooth-part sum.

The density here is the limit of the number of selected ideals of norm at most $x$ divided by $J_K(x)$. We count ideals; associate generators in a ring with infinitely many units would require a separate counting convention.

\clearpage
\section{Where the proofs enter the Lean argument}\label{sec:map}
The readable edition of \texttt{lsernau/grh-lean} at commit \texttt{72fead102f66} has five public declarations. They give Mathlib's RH, its zeta-zero formulation, nonvanishing of any Dirichlet $L$-function in $\Re s>1/2$ apart from the principal pole, the critical-strip statement for all Dirichlet characters, and its primitive specialization. The modulus is any positive integer. The definitions are in \upsrc{src/Statements.lean}{\texttt{Statements.lean}} and the exported theorems in \upsrc{src/Main.lean}{\texttt{Main.lean}}. The project reports a build with Lean 4.35.0, comparator verification against the statements, and kernel exports accepted by independent checkers; the three reported axioms are \texttt{propext}, \texttt{Classical.choice} and \texttt{Quot.sound}. We inspected these source files and the intermediate declarations used here. We did not rebuild the two upstream editions.

\paragraph{The arithmetic estimate.}
Fix a suitable modulus $q$ and auxiliary quadratic character $\chi_d$. The source studies the remainder after subtracting $(\log x)/\varphi(q)$ and a constant from
\[
  \sum_{\substack{n\le x\\n\equiv1\bmod q}}
      \Lambda(n)(1-\chi_d(n))\left(\frac1n-\frac1x\right).
\]
In the \upsrc{src/New/Reduction/MainInequality.lean}{main inequality}, an upper and lower bound apply to the same index-corrected log Gram determinant. Short integral forms supplied by a normalized Poincar\'e map, together with a spanning estimate, give the upper bound; a determinant estimate for a chain of twisted theta lifts gives the lower bound. The main terms cancel, leaving a positive multiple of the arithmetic remainder bounded by the error terms. The chain length is of order $\sqrt{x}$. The \upsrc{src/New/Reduction/OneSidedBound.lean}{one-sided bound} propagates these estimates to real scales. The \upsrc{src/New/Reduction/Zeros/Criterion.lean}{zero criterion} chooses a twist with no zeros in the rectangle under study. Zero multiplicities then enter the Mellin residues with controlled sign, and Landau's theorem excludes off-line zeros. See also Alp\"oge's account \cite{alpoge}.

\paragraph{What carries over.}
The prime test in Section~\ref{sec:binary} preserves that last sign argument, with a group character identity isolating a single Hecke factor. The geometric upper estimate still has to be constructed for the new weight. Two intermediate tools permit unconditional work: \upsrc{src/New/Reduction/Data/Existence.lean}{the sparse-place construction} encodes multiplicative patterns as additive relations, and \upsrc{src/Lib/Cyclotomic/Voronoi.lean}{the Voronoi congruence} governs the local witnesses that bound odd cocycles. The combinatorial and cyclotomic arguments exploit these intermediate tools without the final zero theorem.

\paragraph{Earlier manuscripts.}
The pinned ProveIt directory contains 49 archives. A text comparison covered 268 TeX members from 43 of them. The principal comparisons are the earlier natural-versus-upper-density work \cite{packetgap}, the natural-density and strict-sign results \cite{packetmc}, and the explicit all-witness gcd question \cite{packetcyc}. The accompanying source inventory records the two revisions and the Git blob identifiers used here.
\clearpage
\section{Bounds for sets without rational geometric progressions}\label{sec:density}
A three-term geometric progression of positive integers has the form
$a<b<c$ with $b^2=ac$; its ratio $b/a$ may be a noninteger. Let
$\overline\alpha$ be the supremum of $\limsup_{N\to\infty}|A\cap[1,N]|/N$
over subsets of the positive integers containing no such triple.
For $s\ge2$ write $S_s$ for the positive integers supported on primes at
most $s$, including $1$. Define $\overline\alpha_s$ by forbidding just
those progressions whose ratio, in lowest terms, has numerator and
denominator in $S_s$. The subscripts below are smoothness cutoffs.

\medskip\noindent\textbf{Theorem.}
\[
  0.7365954 < \overline\alpha < 0.7716011,\qquad
  \overline\alpha_7-\overline\alpha_{11}\ge \frac{64}{139755}.
\]
McNew's published bounds are $0.730027<\overline\alpha<0.772059$
(Math. Comp. 84 (2015), 2893--2910, Theorems 4.3 and 4.5;
\href{https://www.nathanmcnew.com/GPFsets.pdf}{author's text}).
The exact supremum remains undetermined. These claims concern upper
density. The incoming report \texttt{mcnew\_density\_and\_grh\_transfers}
separates the two optima for \emph{natural} density and gives restricted
Rankin optima. A natural-density upper bound in that report cannot be
used as an upper bound for $\overline\alpha$.

The source connection is the replacement of multiplicative patterns by
affine relations in exponents. The pinned readable edition applies
\texttt{exists\_three\_term\_free} and \texttt{exists\_large\_isSparse}
in \texttt{New/Reduction/Data/Existence.lean} to pack places with
controlled multiplicative patterns. There it represents places by
powers and constructs exponents using a Behrend sphere. Here $b^2=ac$
gives the coordinatewise relation $2v_p(b)=v_p(a)+v_p(c)$.
McNew already used smooth-number fibers for this density question.
The additions here are explicit finite templates and a strict,
certified loss when the prime $11$ enters the allowed ratios.
The density proofs below are independent of Dirichlet GRH.

\subsection*{Lower bound by an explicit infinite construction}
The file \texttt{lower.csv} gives thresholds
$1=t_1<\cdots<t_{275}=972405$ and sets $C_i\subset S_7\cap[1,t_i]$.
The checker verifies $|C_i|=i$ and tests $b^2/a\notin C_i$
for all $a<b$ in $C_i$ for which $a\mid b^2$. Thus all $C_i$ are
progression-free. They need not be nested. For $x\ge1$ let
$C(x)$ be $C_i$ for the largest $t_i\le x$, and use the empty set
for $x<1$. The exact sum of the threshold contributions is
\[
 I=\sum_{i=1}^{275}\frac1{t_i}
  =\frac{86690601810929969}{26132271943680000}.
\]

For a rough part $h$ require $(h,210)=1$. At primes $11\le p\le997$
allow exponents $0,1,3,4$ and at all larger primes allow $0,1$.
Both exponent sets contain no nonconstant arithmetic progression of
length three. The allowed $h$ have natural density
\[
 c=\frac{6}{\pi^2}\prod_{p\in\{2,3,5,7\}}\frac{p}{p+1}
          \prod_{11\le p\le997}(1+p^{-3}),
\]
where products indexed by $p$ range over primes. Indeed, start with
squarefree integers, impose exponent zero for $p\le7$, and add exponents
$3,4$ for primes in $[11,997]$. At each such prime the local density
factor is multiplied by $1+p^{-3}$. Finite Euler products give the
density by the Chinese remainder theorem and the tail of the
square-divisibility exceptions has upper density at most
$\sum_{p>Y}p^{-2}$, which tends to zero. Here one can bound the
number of exceptions through $N$ by summing $\lfloor N/p^2\rfloor$.

Take integers $N_0=0$, $N_j=2^{2^j}$ for $j\ge1$.
For each allowed $h$ in $(N_{j-1},N_j]$, insert the numbers
$h C(N_j/h)$ into a single set $A$. The rough part of each inserted
number is $h$. If three inserted numbers formed a geometric progression,
their valuations at each prime greater than $7$ would form an
arithmetic progression in a permitted exponent set, so all would have
the same rough part. Their three smooth parts would lie in the one
progression-free template assigned to that part. This proves avoidance
across all stages as well as within a stage. At $N=N_j$, the cardinality
from the current stage alone is bounded below by
\[
  \sum_{i=1}^{275}\#\{h\le N/t_i:\ h\text{ allowed}\}-275N_{j-1}.
\]
This counts the size of each template via threshold indicators. Since
$N_{j-1}/N_j\to0$, division by $N_j$ gives
$\overline d(A)\ge cI$. Replacing $\pi$ by its strict upper bound
$355/113$ yields the rational lower estimate
\[
 \frac{6\cdot113^2}{355^2}
 \prod_{p\in\{2,3,5,7\}}\frac p{p+1}
 \prod_{11\le p\le997}(1+p^{-3})\,I
   >\frac{7365954}{10^7}.
\]
The checker also certifies $\pi<355/113$ using the alternating series in Machin's formula $\pi=16\arctan(1/5)-4\arctan(1/239)$.

\subsection*{A strict loss when the prime eleven is admitted}
Let $F_s(x)$ be the maximum size of a progression-free subset of
$S_s\cap[1,x]$, and let $g_s(N)$ be the maximum size of a subset of
$[1,N]$ avoiding progressions with $s$-smooth ratio.
McNew's Theorem 4.2 gives $g_s(N)/N\to\overline\alpha_s$.
In a fiber with fixed part $b$ coprime to $2310$, put $x=N/b$.
Allowing all $11$-smooth ratios gives the maximum $F_{11}(x)$ in this
fiber. If only $7$-smooth ratios are forbidden, its maximum is
$H(x)=\sum_{j\ge0}F_7(x/11^j)$ because each $11$-adic layer can be
chosen independently. Always $H(x)\ge F_{11}(x)$.
The following finite comparison gives an additional exclusion throughout
$121\le x<165$.
\[
\begin{array}{c|rrrrrrr}
 [a,b') &[121,125)&[125,132)&[132,135)&[135,140)&[140,154)&[154,162)&[162,165)\\
 \hline
 \max(S_{11}\cap[1,b')) &121&128&132&135&150&160&162\\
 H(a)\ \text{at least} &38&39&40&41&42&43&44\\
 F_{11}(x)\ \text{at most} &37&38&39&40&41&42&43
\end{array}
\]
For each interval, \texttt{intervals.json} supplies three $7$-smooth
sets for layers $j=0,1,2$. Testing their members and their triples
proves the lower entry for $H(a)$. For an upper entry $u$, use the
hypergraph whose vertices are the $11$-smooth numbers up to the listed
maximum, with a hyperedge for every triple $d<e<f$ satisfying $df=e^2$.
A set of $u+1$ surviving vertices would give a cover with budget
$|V|-u-1$ deleted vertices. The seven proof trees in \texttt{trees/}
rule this out by the same exact checks.

At an internal node, branch on a single vertex being deleted or kept.
Discard the covered edges and let $k$ be the remaining deletion budget.
A leaf closes if $k<0$ or a still uncovered edge has all its vertices
fixed to be kept. The remaining type of leaf carries positive integer
weights on uncovered edges. Let $W$ be their sum and $D$ the maximum
total weight incident to any undecided vertex. Every cover using at most
$k$ more deletions covers weight at most $kD$. The supplied weights
satisfy $W>kD$. The verifier checks these integer inequalities, the
two branches at each internal node, all listed members and the
ordering of the smooth numbers. The seven trees contain respectively
$813,3295,1559,1331,2047,3139,2865$ nodes. Their validity is independent
of the numerical optimizer used to find leaf weights.

Fibers with $(b,2310)=1$ partition $[1,N]$ for both optimization
problems. Every $b$ with $N/165<b\le N/121$ costs at least one in the
$11$-smooth problem. Dividing $g_7(N)-g_{11}(N)$ by $N$ and passing
to the limit therefore gives
\[
 \overline\alpha_7-\overline\alpha_{11}
 \ge \frac{\varphi(2310)}{2310}
       \left(\frac1{121}-\frac1{165}\right)=\frac{64}{139755}.
\]
McNew's section 4 obtains the more specific bound
$\overline\alpha_7<0.772059$ immediately before Theorem 4.3.
Consequently
\[
 \overline\alpha\le \overline\alpha_{11}
 < \frac{772059}{10^6}-\frac{64}{139755}
 =\frac{21567021109}{27951000000}<0.7716011.
\]

\clearpage
\section{The exact gcd of the cyclotomic annihilator witnesses}
\label{sec:exact-stickelberger}

The 26 November ProveIt manuscript asks for the gcd of all positive integers
satisfying both clauses of the Lean predicate
\texttt{New.OddPart.IsStickelbergerBound}. Its theorem determines the gcd
for a family constructed by clearing denominators before applying Voronoi's
congruence. Clearing those denominators later makes the full problem tractable.
This section gives the exact gcd at every modulus and weight, a finite
decision procedure for the least individual witness, and arithmetic consequences
for the cocycles in the repository.
The earlier question and its local auxiliary-prime calculation are in
\href{https://github.com/VladimirReshetnikov/ProveIt/blob/a1d43d8a32add57d596273e04648bc96a0a9bc0d/docs/incoming/grh_method_transfers_2026-11-26.zip}{the prior manuscript, section on the cyclotomic annihilator}.

\subsection{The predicate and its exact character denominator}

Fix integers $N,e\ge1$. A Dirichlet character $\chi$ modulo $N$ is relevant
when $\chi(-1)=(-1)^{e+1}$. Characters take their usual value zero on
nonunits. For $N=1$ the unique character takes value one on every integer.
Let $O_\chi=\mathbb Z[\chi((\mathbb Z/N\mathbb Z)^\times)]$.
With $B_j(x)$ the Bernoulli polynomial set
\begin{align*}
 K&=(e+2)!(e+1)N,&
 b(a)&=(e+2)!N^{e+1}B_{e+1}(a/N),\\
 B_\chi&=\sum_{a=0}^{N-1} b(a)\chi(a),&
 F(c)&=(c^e-1)(c^{e+1}-1)\quad(c\ge2).
\end{align*}
Upstream shows $b(a)\in\mathbb Z$ and $B_\chi=-K L(-e,\chi)\ne0$ for
relevant $\chi$, including imprimitive characters with their actual Euler factors.
For $m\ge1$ and $c\equiv1\pmod N$ put
\[
 S_{m,c,\chi}=\sum_{a=0}^{Nm-1}
   \left\lfloor\frac{ca}{Nm}\right\rfloor a^e\chi(a).
\]
A positive integer $X$ satisfies the predicate when both conditions hold.
First, for every prime $p$, some $c\ge2$ congruent to $1$ mod $N$ has
$p\nmid c$ and $F(c)\mid X$. Second, for each $c\ge0$ congruent to $1$
mod $N$ and each relevant character in a characteristic-zero integral domain
$O$, there is a single $w\in O$ such that
\[
  (c^{e+1}-1)X\equiv 2w S_{m,c,\chi}\pmod{mO}
  \qquad\text{for every }m\ge1\text{ with }(m,c)=1.             \tag{1}
\]
The upstream Voronoi congruence, in the character ring, reads
\[
 Kc^e S_{m,c,\chi}\equiv (c^{e+1}-1)B_\chi\pmod{mO_\chi}.    \tag{2}
\]
The interface is stated in lines 136--151, the special-value identities
in lines 153--203, and (2) in lines 245--256 of
\href{https://github.com/lsernau/grh-lean/blob/72fead102f663d5c7a407964a13dc0cf9d977207/grh/src/Lib/Cyclotomic/Voronoi.lean}{\texttt{Voronoi.lean}}.

\paragraph{Theorem 1.}
The second clause holds precisely when
\[
       \frac{KX}{2B_\chi}=-\frac{X}{2L(-e,\chi)}\in O_\chi
       \quad\text{for each relevant complex character }\chi.       \tag{3}
\]
In particular, it places exactly the condition $Q\mid X$ on a positive
integer $X$, where $Q$ is the least positive integer clearing the
denominators of $(2L(-e,\chi))^{-1}$ in all the $O_\chi$. Set $Q=1$ if
there are no relevant characters.

\paragraph{Proof of necessity.}
Apply the second clause to the character in $O_\chi$ itself. Fix $c>1$,
multiply (1) by $Kc^e$, and use (2). Then
$(c^{e+1}-1)(Kc^eX-2wB_\chi)$ belongs to $mO_\chi$ for all $m$ coprime
to $c$. The intersection of these ideals is zero, since $O_\chi$ is a
finite free abelian group and we may take powers of a single prime outside
$c$. Hence $w=c^e KX/(2B_\chi)$. Use $c=1+N$ and $c=1+2N$.
Their $e$th powers are coprime, so B\'ezout's identity gives (3).

\paragraph{Proof of sufficiency.}
Write $R=Nm$. For $k\ge1$ coprime to $c$, partition the terms of $S_{km,c,\chi}$
as $a+jR$, $0\le a<R$, $0\le j<k$. Modulo $m$ the polynomial and character
factors are unchanged. The floor factors satisfy
\[
 \sum_{j=0}^{k-1}\left\lfloor\frac{ca/R+cj}{k}\right\rfloor
  =\left\lfloor\frac{ca}{R}\right\rfloor+\frac{(c-1)(k-1)}{2}.
\]
Indeed the residues of $cj+\lfloor ca/R\rfloor$ modulo $k$ run once
through $0,\ldots,k-1$. Let $A_m=\sum_{a=0}^{R-1}a^e\chi(a)$.
The involution $a\mapsto -a\bmod R$ negates its summands modulo $m$,
using the parity of $\chi$; the zero summand causes no exception since
$e\ge1$. Thus $2A_m\equiv0\pmod m$, even when $m$ is even.
The floor identity, whose correction is an integer, gives
\[
                2S_{km,c,\chi}\equiv2S_{m,c,\chi}\pmod m.          \tag{4}
\]
Take $w=c^eKX/(2B_\chi)\in O_\chi$ and denote the left side of (1) minus
its right side by $E_m$. For any $m$ prime to $c>0$ put
$k=\prod_{p\mid m}p^{v_p(K)}$. Then $(km,c)=1$. At the enlarged modulus
$km$, (2) and $2w B_\chi=Kc^eX$ give $Kc^e E_{km}\in kmO_\chi$.
Cancel $k$ in the integral domain and invert $(K/k)c^e$ modulo $m$.
It follows that $E_{km}\in mO_\chi$. Equation (4) now gives
$E_m\in mO_\chi$. For $c=0$, possible only at $N=1$, the sole modulus
coprime to $c$ is $m=1$. The argument proves the result in all target
domains $O$ by passage along the character-ring embedding. Indeed, the finite
nonzero image in $\operatorname{Frac}(O)^\times$ is cyclic of order $d$, and
the image of $\mathbb Z[\zeta_d]$ under a choice of primitive generator is
the character subring. Choosing the corresponding complex character
transports every congruence. This finishes the proof.

\subsection{All witnesses and the least witness}

Let $D$ be the least positive integer divisible by every $2B_\chi$
in the corresponding $O_\chi$, as in the earlier report. By definition
$D\mid KX$ is equivalent to (3), so
\[
                           Q=D/\gcd(D,K).
\]
For each prime $p$ let $m_p$ be the minimum of $v_p(F(c))$ over
$c\ge2$, $c\equiv1\pmod N$, $p\nmid c$. The first-clause computation of
the earlier report gives, for odd $p$,
\[
 m_p=\begin{cases}
  2v_p(N)+v_p(e(e+1)),&p\mid N,\\
  \displaystyle\sum_{j\in\{e,e+1\}:\;p-1\mid j}(1+v_p(j)),&p\nmid N.
 \end{cases}
\]
At $p=2$ it is $v_2(e(e+1))+2v_2(N)$ for $4\mid N$ and
$v_2(e(e+1))+3$ otherwise. For odd $p\mid N$, the lifting-the-exponent
formula gives the minimum when $v_p(c-1)=v_p(N)$.
For odd $p\nmid N$, choose a primitive root modulo $p^2$ in the permitted
progression; Fermat's theorem and lifting give the lower bound.
For odd $c$, the total at two is
$2v_2(c-1)+v_2(c+1)-1+v_2(e(e+1))$. When $4\nmid N$, take
$c\equiv3\pmod8$ to attain the minimum. Write $M=\prod_p p^{m_p}$,
a finite product.

\paragraph{Theorem 2.}
The gcd of \emph{all} positive witnesses of the upstream predicate is
\[
              G(N,e)=\operatorname{lcm}(Q,M).                     \tag{5}
\]
Moreover, a specified $X>0$ is a witness precisely when $Q\mid X$
and the finite set $\{c\ge2:c\equiv1\pmod N,\ F(c)\mid X\}$ has
gcd one (with empty gcd zero).

\paragraph{Proof.}
The first clause says precisely that the eligible $c$ have no common
prime divisor and are nonempty. It implies $M\mid X$; Theorem~1 gives
$Q\mid X$. Conversely, for coprime $a,b\ge2$ in the progression,
$X(a,b)=\operatorname{lcm}(Q,F(a),F(b))$ is a witness.
For each prime $p$ choose a minimizer $a$ for $m_p$. By the Chinese
remainder theorem choose $b>1$ congruent to $a$ modulo $p^{m_p+1}$, to
$1$ modulo $N$, and to $1$ modulo every prime divisor of $a$. These
congruences are compatible since $p\nmid a$ and $a\equiv1\pmod N$.
Hence $(a,b)=1$ and $v_p(F(b))=m_p$. The $p$-valuation of $X(a,b)$
is $\max(v_p(Q),m_p)$, as required. Finitely many witnesses suffice to
realize the gcd: start with any witness and choose the corresponding
minimizer for each of its prime divisors.

The least witness is a different invariant. Start with the explicit upper
bound $U=X(1+N,1+2N)$ and list all $c$ in the progression with $F(c)\le U$.
For each subset $T$ of the list with $\gcd(T)=1$ compute
$\operatorname{lcm}(Q,\{F(c):c\in T\})$. The smallest value is the least
witness, by the finite characterization just proved. The accompanying
code merges equal pairs of partial gcd and partial lcm during this search.

\subsection{Numerical examples and the arithmetic consequence}

\begin{center}
\small
\begin{tabular}{r r r r r r}
\hline
$N$ & $e$ & $Q$ & $M$ & previous $\operatorname{lcm}(D,M)$ & exact $G$\\
\hline
1&11&691&131040&286858091520&90548640\\
3&9&9841&14256&3928212288000&140293296\\
5&1&4&1200&1200&1200\\
5&2&4&1200&7200&1200\\
7&2&32&2352&112896&4704\\
8&2&6&384&3456&384\\
12&2&20&864&17280&4320\\
15&1&16&3600&14400&3600\\
\hline
\end{tabular}
\end{center}

For $N=1,e=11$ the special value is $\zeta(-11)=691/32760$, giving
$Q=691$ immediately. The prime $11$ divides the previous bound and is
absent from $G$. At $(N,e)=(3,9)$ only the principal character is relevant:
$L(-9,\chi_0)=(1-3^9)\zeta(-9)=9841/66$. Here both $5$ and $7$
are removed. At $(8,2)$ the least individual witness is $48\,746\,880$,
and its eligible auxiliaries are $9,25$. The pair $9,17$ would instead
give $643\,668\,480$. These minima and the table can be reproduced from
the exact polynomial and integer operations of
\texttt{check\_stickelberger.py}. Its output also gives finite lists of
witnesses realizing each gcd, and tests (1) directly on 1,217
character, auxiliary and modulus choices across eleven examples.

Fix a valid $X$. The source proves that $4X^2$ kills the cocycles in $\mathsf{component}(\ell,e,\psi,\nu)$. These are locally constant cocycles with quotient coefficients for the inverse of $\mathrm{cyc}^e\psi$, subject to the stated inertia and local vanishing conditions, with $\mathrm{cyc}^e\psi$ odd. Taking integer linear
combinations over a finite family realizing $G$ proves that $4G^2$
kills it. Thus at $(N,e)=(3,9)$ all of these cocycles vanish at each level
for coefficient prime $5$ or $7$: multiplication by $4G^2$ is a unit
on a module with coefficients modulo the corresponding prime power.
The same reasoning improves the source's finite-level CM descent
annihilator to $16\varphi(N)G^2$ under its own hypotheses. Optimality
in (5) concerns the stated predicate. A module can have a smaller
annihilator, and a surviving prime may still have zero component.
Classical higher Stickelberger theory treats related Galois-cohomological
annihilators; the question solved here comes from the specific Lean
congruence and its preceding analysis.
See
\href{https://github.com/lsernau/grh-lean/blob/72fead102f663d5c7a407964a13dc0cf9d977207/grh/src/Lib/Cyclotomic/Kolyvagin/ReciprocityComponentBounds.lean\#L670-L680}{the component annihilation lemma},
\href{https://github.com/lsernau/grh-lean/blob/72fead102f663d5c7a407964a13dc0cf9d977207/grh/src/Lib/Cyclotomic/OddSha.lean\#L529-L545}{the CM descent lemma},
and \href{https://doi.org/10.1016/j.crma.2003.09.019}{Banaszak, \emph{Higher analogues of Stickelberger's theorem}}.

\clearpage
\section{How many sparse places can the construction use?}
\label{sec:spcap}

The choice of infinite places in the Lean proof poses a finite extremal
problem on residue classes. Two of its prime constraints admit an exact
joint solution. Keeping the other prime constraints gives both an upper
bound and a larger explicit construction.

\paragraph{The finite problem.}
Fix a positive multiple $q_0$ of $18480$, and let $s\ge1$ be an integer.
Put
\[
 z=3^{\lfloor\log_3s\rfloor+4},\qquad q_1=q_0\,3z\,2^s,\qquad
 C=\{x\in(\mathbb Z/q_1\mathbb Z)^\times:x\equiv-1\pmod{q_0}\}.
\]
For a prime $p$, let $m_p=q_1/p^{v_p(q_1)}$. Say that $T\subset C$
satisfies the condition at $p$ if there are no $x,y,t\in T$ and
$0<i<j<z$ with
\[
  y/x\equiv\pm p^i\pmod {m_p},\qquad
  t/x\equiv\pm p^j\pmod {m_p}.
\]
The signs are independent and the three places may coincide. Imposing
this condition at all primes $p\le L$ is exactly the sparse-set requirement
on the $b$-places. The source gives the
\href{https://github.com/lsernau/grh-lean/blob/72fead102f663d5c7a407964a13dc0cf9d977207/grh/src/New/Reduction/Data/FieldsAndPlaces.lean\#L65-L110}{moduli and tower parameters}
and the
\href{https://github.com/lsernau/grh-lean/blob/72fead102f663d5c7a407964a13dc0cf9d977207/grh/src/New/Reduction/Data/AdmissibleData.lean\#L54-L62}{definition of sparsity}.
Its
\href{https://github.com/lsernau/grh-lean/blob/72fead102f663d5c7a407964a13dc0cf9d977207/grh/src/New/Reduction/Data/Existence.lean\#L164-L238}{construction of sparse places}
fixes one odd residue and uses a large set of exponents in the $2$-primary
coordinate.

\paragraph{Exact two-prime theorem.}
For $p=2,3$ write $b_p=q_0/p^{v_p(q_0)}$ and $N_p=m_p/b_p$, so
$N_2=3z$ and $N_3=2^s$. Let $d_p$ be the least positive integer for
which $p^{d_p}\equiv\pm1\pmod{b_p}$. Let $g_p=\epsilon_p p^{d_p}$,
where $\epsilon_p\in\{-1,1\}$ is the unique sign giving
$g_p\equiv1\pmod{b_p}$. Define
\[
 H_p=\operatorname{ord}_{m_p}(g_p),\qquad
 w_p=\lceil z/d_p\rceil,\qquad
 A_p=\frac{N_p}{H_p}\min\left(H_p,
           \left\lfloor\frac{2H_p}{w_p}\right\rfloor\right).
\]
The largest size of a set $T\subset C$ satisfying the conditions at both
$2$ and $3$ is $A_2A_3$. Thus $A_2A_3$ is the exact optimum for $L=3,4$,
and an upper bound for every $L\ge3$. In particular,
\[
    \frac{|T|}{|C|}\le \frac{4d_2d_3}{z^2}\quad (L\ge3).
\]
Both $d_p$ depend only on the fixed modulus $q_0$.

\paragraph{Proof.}
The possible images in $(\mathbb Z/m_p\mathbb Z)^\times$ are precisely
the residues $a\equiv-1\pmod{b_p}$. They number $N_p$, since every
prime divisor of $m_p$ already divides $b_p$. If $\pm p^i$ takes one
of these residues to another, then $d_p\mid i$ by the order of the
class of $p$ in $(\mathbb Z/b_p\mathbb Z)^\times/\{\pm1\}$. Since
$b_p>2$, the sign is unique and the step is multiplication by
$g_p^{i/d_p}$. Multiplication by $g_p$ partitions the possible images
into $N_p/H_p$ cycles of length $H_p$.

Write $H=H_p$ and $w=w_p$. On one such cycle mark the residues in the
image of $T$, reading their indicator as a periodic binary sequence on
$\mathbb Z$. Three marked positions in a window of $w$ consecutive
indices give a forbidden triple by subtracting the first index and
multiplying both differences by $d_p$. Conversely, the three indices in any
forbidden triple fit in such a window because
$\lfloor(z-1)/d_p\rfloor=w-1$. If a window goes around the cycle twice,
all its occurrences count, as required for repeated places. For $K$ marks
the sum of occupancies of the $H$ windows is $wK$, so $K\le H$ and
$K\le\lfloor2H/w\rfloor$. Both bounds are attained together. With
$K=\min(H,\lfloor2H/w\rfloor)$ mark the indices with
\[
 \left\lfloor\frac{(h+1)K}{H}\right\rfloor-
 \left\lfloor\frac{hK}{H}\right\rfloor=1 .
\]
There are $K$ such indices per period, and the sum on every window is at
most $\lceil wK/H\rceil\le2$ by telescoping. This proves the formula for
the largest allowable projected support $A_p$.

Write $q_0=2^v3^a b$ with $\gcd(b,6)=1$. The component modulo $b$ is
fixed to $-1$ in every element of $C$. The two remaining freely varying
components are a lift of $-1\bmod 2^v$ and a lift of $-1\bmod 3^a$.
The condition at $2$ uses only the second of these, and the condition
at $3$ only the first. Any two permissible projected supports can
therefore be combined by taking their full Cartesian product under
the Chinese remainder theorem. Every eligible $T$ is contained in the
product of its own projections. This proves exactness. Finally
$A_p/N_p\le 2/w_p\le2d_p/z$ and $|C|=N_2N_3$.\hfill$\square$

\paragraph{The base modulus $18480$.}
Here $b_2=1155$ and $b_3=6160$. The orders of $2$ modulo $3,5,7,11$
are $2,4,3,10$. The powers of $2$ cannot reach $-1$ modulo $7$. The orders of $3$
modulo $16,5,7,11$ are $4,4,6,5$. The powers of $3$ cannot reach $-1$ modulo $16$.
Consequently $d_2=d_3=60$ and both signs are positive. The identities
\[
 v_3(2^{60}-1)=2,\qquad v_2(3^{60}-1)=4
\]
give $H_2=z$, $H_3=2^s$. Indeed, in the first case replacing the
exponent $60$ by $60k$ adds $v_3(k)$ to the valuation. In the second
case it adds $v_2(k)$, since $3^{60}\equiv1\pmod4$. These elementary
lifting identities follow by factoring powers with exponent prime to
the valuation prime and then expanding the next prime power by the
binomial theorem. The other prime-power components already have order
one. With $w=\lceil z/60\rceil$, the answer is therefore
\[
 3\min\bigl(z,\lfloor 2z/w\rfloor\bigr)\,
   \min\bigl(2^s,\lfloor 2^{s+1}/w\rfloor\bigr).
\]
For $s\ge81$ this simplifies to
$357\lfloor2^{s+1}/\lceil z/60\rceil\rfloor$.
In fact $z=6561$ gives $w=110$ and $\lfloor2z/w\rfloor=119$.
The next possible $z$ is $19683$. For it and all larger $z$,
$120z/(z+59)>119$, while $2z/w<120$ because $60\nmid z$.
The possible $3z$ odd residues have thus shrunk to exactly $357$ before
any of the primes $p\ge5$ are considered.

\paragraph{An all-prime construction.}
The exact extremizer at $2$ can also be retained when imposing \emph{every}
prime up to $L$. Let $q_0=2^v u$ with $u$ odd. For positive integers $n,h$
assume
\[
 ((z-1)h)^n\le2^s,\qquad (L+1)(z-1)<2^{v+s}.
\]
There exists $T\subset C$ satisfying all conditions at $p\le L$ with
\[
                    |T|\ \ge\ A_2\frac{h^n}{n h^2}.
\]
For the same $n,h$ this multiplies the cited lower bound by $A_2$
and relaxes its two sufficient size conditions.

To prove it, choose $S\subset[0,2^s)\cap\mathbb Z$ of size at least
$h^n/(nh^2)$ such that every weighted relation
$j e_2\equiv (j-i)e_1+i e_3\pmod{2^s}$ with $0<i<j<z$ forces
$e_1=e_2=e_3$. This is the source's
\href{https://github.com/lsernau/grh-lean/blob/72fead102f663d5c7a407964a13dc0cf9d977207/grh/src/New/Reduction/Data/Existence.lean\#L77-L117}{weighted Behrend lemma}
with coefficient budget $z-1$. In brief, encode the vectors of
$\{0,\ldots,h-1\}^n$ using base $(z-1)h$. At most $nh^2$ distinct
squared lengths occur, so some sphere contains at least $h^n/(nh^2)$
vectors. In a weighted relation all digit sums
are below the base, so each integer is strictly less than $((z-1)h)^n\le2^s$. Equality modulo $2^s$ is therefore equality coordinatewise. Strict
convexity of the squared norm forces equality on the sphere.

Take as the $2$-primary component $-(1+2^v)^e\bmod 2^{v+s}$ for $e\in S$,
and independently take an optimal support of size $A_2$ in the odd
modulus $m_2$. The first components are distinct since $1+2^v$ has
order $2^s$, by the same binomial lifting identity. The condition at $2$
holds by construction. Let $p\le L$ be odd and let $t_p$ be the choice
of $p$ or $-p$ that equals $1\bmod4$. Project a proposed forbidden
triple to $2^{v+s}$, a divisor of $m_p$. The two ratios are $1\bmod4$,
so their signs are forced and their values are $t_p^i,t_p^j$.
Raising the first to the $j$th power and the second to the $i$th
gives the weighted relation on exponents. All three exponents are equal,
which would give $t_p^i\equiv1\pmod{2^{v+s}}$. Yet
\[
 2^{v_2(t_p^i-1)}=2^{v_2(t_p-1)}2^{v_2(i)}
        \le(p+1)i<2^{v+s},
\]
a contradiction. The argument explains why the coefficient budget can
drop from $2z$ to $z-1$: the sign is determined modulo $4$ before the
powers are compared.\hfill$\square$

\paragraph{Sharp order for every fixed prime cutoff.}
Let $M(q_0,s,L)$ denote the maximum over sets satisfying every condition
at $p\le L$. For each fixed $q_0$ as above and each fixed $L\ge3$,
\[
  M(q_0,s,L)=\Theta_{q_0,L}(2^s/z),\qquad
  \frac{M(q_0,s,L)}{|C|}=\Theta_{q_0,L}(z^{-2})
  \quad(s\longrightarrow\infty).
\]
The upper bound is the two-prime theorem. For the lower bound, form a
graph on the $2^s$ possible $2$-primary components. Join $x$ to $y$
if $y/x$ is a signed $p^i$ for some odd prime $p\le L$ and $0<i<z$,
or if $x/y$ has this form. The preceding valuation bound rules out
loops once $(L+1)(z-1)<2^{v+s}$, which holds for all large $s$ since
$z\le81s$. For each pair $p,i$, the residue modulo $4$ forces the sign,
so the degree is at most $2r(z-1)$ where $r$ counts the odd primes up to $L$.
Greedily choosing a vertex and discarding its neighbors yields an
independent set of size at least $2^s/(1+2r(z-1))$. This excludes even
the first nontrivial step of each forbidden triple at an odd prime.
Use its full product with an optimal even-prime support.
The latter has $A_2\ge1$: a singleton odd residue could return under
a signed $2^i$ only if $4^i\equiv1\pmod{3z}$, whereas $4=1+3$
has order $z$ there. Consequently
\[
 M(q_0,s,L)\ \ge\ \frac{A_2\,2^s}{1+2r(z-1)}.
\]
This proves the claimed order and gives an explicit lower bound whenever
the displayed size condition on $L$ holds.\hfill$\square$

\paragraph{Reproduction.}
The companion program \path{sparse_places.py} independently evaluates
the two orders, enumerates all binary cycle supports for $108$ small
size/window pairs, and tests explicit attaining supports against the
signed-power predicate. For $q_0=18480$, $s=40$, $L=19$, $n=2$, $h=40$,
it takes eight exponents on a single digit sphere and $354$ projected
residues at $2$, obtaining $2832$ eligible places. The source's
construction and the two extremal formulas are finite combinatorics
independent of RH. Determining the exact all-prime capacity for the
growing cutoff used in the Lean proof remains a separate question.

\clearpage
\section{Which binary Frobenius tests can discard selected factors?}
\label{sec:binary}
At an unramified prime one often records a yes/no event, such as whether a polynomial stays irreducible. A zero of an $L$-function contributes to the explicit formula with a coefficient determined by that event. For a one-sided estimate to detect just a selected family, one wants all its coefficients to have the same sign and the unwanted coefficients to vanish. The earlier ProveIt report \cite{packetmc} classifies events with strictly negative coefficients on \emph{every} nonlinear irreducible. Here we allow zero coefficients. The possibility of equality changes the answer: a quotient, a normal subgroup in that quotient, and a sign character can all occur.

\subsection{Classification}
Let $H\triangleleft G$ be a normal subgroup of a finite group, $Q=G/H$, and let $f:G\to\{0,1\}$ be constant on conjugacy classes with $f(x^{-1})=f(x)$. On $\mathbb R^G$ write
\[
 A_{xy}=f(x^{-1}y),\qquad
 V=\{v:\ \sum_{x\in C}v_x=0\ \text{for every }H\text{-coset }C\}.
\]
The orthogonal projection onto coset-constant vectors averages over $H$ and commutes with $A$. For an irreducible character $\chi$, $A$ acts on its isotypic space by $|G|\langle f,\chi\rangle/\chi(1)$, with the normalized character inner product. The characters in $V_{\mathbb C}$ are precisely those nontrivial on $H$. In particular $H=[G,G]$ singles out all the nonlinear characters. Call $f$ admissible if $A|_V\le0$.

\begin{theorem}[Classification with zero coefficients]\label{th:binary}
Coset-constant events are admissible and annihilate $V$. Every other admissible event factors through a quotient $G'=G/K$ for some $K\triangleleft G$ strictly contained in $H$. Write $H'=H/K$, $q:G'\to Q$. There exist a normal subgroup $J\triangleleft G'$ with $J\cap H'=1$, a conjugation-invariant homomorphism $\varepsilon:J\to\{1,-1\}$, and a binary class function $b$ on $Q$, invariant under inversion, such that
\begin{equation}\label{eq:binary}
 f(x)=b(q(x))-\varepsilon(x)\mathbf 1_J(x)\qquad(x\in G').
\end{equation}
Here $\varepsilon\mathbf 1_J$ is extended by zero and $b(q(j))=(1+\varepsilon(j))/2$ for $j\in J$. Conversely every such event is admissible. On $V$ its only possible nonzero eigenvalue is $-|K||J|$, with multiplicity
\[
               \frac{|G/K|-|G/H|}{|J|}.
\]
For a character nontrivial on $H$, a coefficient survives exactly when $K$ acts trivially and $J$ acts on the resulting representation of $G'$ by the scalar character $\varepsilon$.
\end{theorem}

\begin{proof}
For $u,v$ in the same fiber, the negative of the quadratic form on $e_u-e_v$ is $2(f(u^{-1}v)-f(1))$. If it is zero, then $A(e_u-e_v)=0$, by semidefiniteness and invariance of $V$. If $f(1)=1$, all these differences give zero, so $f$ is constant in each fiber. Otherwise $f(1)=0$. The elements $h\in H$ for which $f(h)=0$ are exactly those with $A e_h=A e_1$. By translation invariance they are the intersection of $H$ with the stabilizer of a column, hence form a subgroup $K$; centrality of $f$ makes $K$ normal. Column equality and translation show $f$ is constant on $K$-cosets. If $K=H$ the first case applies. Otherwise pass to $G/K$. The operator on the $K$-coset-constant subspace becomes $|K|$ times the quotient operator and vanishes on its orthogonal complement. We now work in that quotient and write $h=|H'|\ge2$.

In the identity fiber $f$ is one off the identity and zero at the identity. Split $V$ into spaces $W_C$ of sum-zero vectors on individual $H'$-cosets. For $B=-A|_V$ every diagonal block is the identity on $W_C$. Each off-diagonal block is a contraction, by positivity of the two-by-two block compression with identity diagonal. For fibers $C,D$, the binary matrix of $A$ between them is regular of row and column sum $r=\sum_{x\in C^{-1}D}f(x)$. Its squared Frobenius norm on all vectors is $hr$, and the constant direction accounts for $r^2$. As at most $h-1$ remaining singular values are nonzero and each is at most one, $r(h-r)\le h-1$. Hence $r$ is $0,1,h-1$ or $h$.

A constant fiber gives a zero block. Any other fiber of $q$, called active, has one exceptional element $s_a$ over $a\in Q$: a missing element among ones, or a present element among zeros. Give it sign $\varepsilon_a=+1$ in the first case and $-1$ in the second. When $h=2$ always choose the missing element. On $W$ the corresponding block of $B$ is $\varepsilon_a$ times the permutation identifying the two fibers by right multiplication by $s_a$. The identity fiber is active, with $s_1=1$ and positive sign.

Write the positive block matrix as a Gram matrix $B_{CD}=U_C^*U_D$, where each $U_C$ is an isometry. A block that is an isometry identifies equal image subspaces, since orthogonal projection onto one image preserves the norms of all vectors of the other and the dimensions agree. Consequently if $a,b$ are active, the blocks between the cosets $1,a,ab$ give
\[
                 B_{1,ab}=B_{1,a}B_{a,ab}.
\]
It is an isometry, so $ab$ is active. The active cosets therefore form a subgroup of $Q$. In the last identity the product of the exceptional matchings is right multiplication by $s_a s_b$. On the sum-zero space a permutation is determined by its restriction. For $h\ge3$ no permutation is minus the identity on that space: it would have to interchange each pair of distinct coordinates. Thus $s_{ab}=s_a s_b$ and $\varepsilon_{ab}=\varepsilon_a\varepsilon_b$. For $h=2$ all chosen signs are positive and ordinary faithfulness gives the same conclusions. Conjugation preserves the chosen exceptional element and sign in every active fiber. Their elements therefore form a normal subgroup $J$ with $J\cap H'=1$, and their signs give an invariant character. Filling in the constant background in each fiber gives \eqref{eq:binary}.

Conversely the first term of \eqref{eq:binary} annihilates $V$. On the $K$-invariant part of $V$, the operator $-A$ acts as $|K|\sum_{j\in J}\varepsilon(j)R_j$, where $(R_j v)(x)=v(xj)$. The latter is $|K||J|$ times the orthogonal projection to the $\varepsilon$-isotypic subspace for $J$. On all $K$-invariant vectors this projection has rank $|G/K|/|J|$ by the regular action of $J$ on each of its cosets. On $H$-invariant vectors the rank is $|G/H|/|J|$, since $J$ embeds in $Q$. Subtraction proves the rank formula. Invariance of $\varepsilon$ and normality of $J$ make this projection central, which gives the criterion on irreducibles.
\end{proof}

\paragraph{An obstruction.}
For $G=S_n$, $n\ge5$, and $H=A_n$, a binary event with all nonlinear coefficients nonpositive either kills them all or detects them all. Indeed $A_n$ is simple, so in the nonconstant case $K=1$. Normal subgroups with trivial intersection commute, hence $J\subset C_{S_n}(A_n)=1$. To see the last equality, an element centralizing every 3-cycle preserves the support of every 3-cycle, hence every point. The theorem now gives the same eigenvalue $-1$ on all of $V$. This rules out selective binary tests in these groups under the common-sign condition.

\subsection{A minimal event for a twisted cubic}
Let $L$ be the splitting field of $P(X)=X^3-X-1$ and set $E=L(i)$. Reduction modulo $2$ is irreducible and the discriminant is $-23$. Thus $\mathrm{Gal}(L/\mathbb Q)=S_3$ and its unique quadratic field is $\mathbb Q(\sqrt{-23})$. It follows that $G=\mathrm{Gal}(E/\mathbb Q)=S_3\times C_2$. Let $\sigma$ be the sign character, $\rho$ the two-dimensional standard character, and $\eta$ the nontrivial character of $C_2$. The untwisted example and its associated weight-one form are described in \cite{serrepoly,ng}.

For a prime $p\nmid46$ define $f(p)=1$ in either of the following cases, and $f(p)=0$ otherwise.
\begin{center}
\begin{tabular}{@{}ll@{}}\toprule
Congruence of $p$ & Factorization of $P$ modulo $p$ \\\midrule
$p\equiv1\pmod4$ & irreducible \\
$p\equiv3\pmod4$ & three distinct linear factors \\\bottomrule
\end{tabular}
\end{center}
For the corresponding class function one has
\begin{equation}\label{eq:identity}
          f=\frac{1+\sigma}{4}+\frac{(1+\sigma)\eta}{12}
                      -\frac{\rho\eta}{3}.
\end{equation}
This follows by evaluating on the three cycle types for each value of $\eta$; $\rho$ has traces $2,0,-1$ on identity, transpositions and 3-cycles respectively. The mean is $1/4$. There is no $\rho$ term.

\begin{corollary}[Smallest support]
This event is the unique least-support inversion-invariant binary class function having negative coefficient on $\rho\eta$ and zero coefficient on $\rho$. Its support has three of the twelve elements.
\end{corollary}
\begin{proof}
Here $H=A_3\times1$ has order three. For a selective test in Theorem~\ref{th:binary}, $K=1$ and $J\ne1$. Since $J$ centralizes $H$, it lies in $C_3\times C_2$. Disjointness from $H$ leaves the unique subgroup $1\times C_2$. To select $\rho\eta$, the sign must be $\eta$. The active fibers therefore have two selected elements and one selected element, respectively; choosing zero on the remaining two fibers is the unique minimum.
\end{proof}

\begin{theorem}[One-sided criterion for one Hecke $L$-function]\label{th:cubic}
Assume Dirichlet GRH as supplied by the upstream theorem. All nontrivial zeros of $L(s,\rho\eta)$ lie on $\Re s=1/2$ if and only if for every $\epsilon>0$ there is a constant $C_\epsilon$ such that, for $x\ge2$,
\begin{equation}\label{eq:cubic}
       \sum_{p\le x,\ p\nmid46} f(p)\log p-\frac{x}{4}
                   \le C_\epsilon x^{1/2+\epsilon}.
\end{equation}
\end{theorem}
\begin{proof}
The standard representation of $S_3$ is induced from either nontrivial character of $A_3$. Twisting by $\eta$ commutes with induction. By Artin formalism and abelian class field theory, $L(s,\rho\eta)$ is the Hecke $L$-function of a nontrivial finite-order character over $\mathbb Q(\sqrt{-23})$, hence is entire, with the usual functional equation. The representation is self-dual. Only this monomial case of Artin holomorphy is needed.

Write $g_p$ for Frobenius and $\Psi_f(x)=\sum_{p^k\le x,\ p\nmid46} f(g_p^k)\log p$. Prime powers with $k\ge2$ contribute $O(\sqrt x\log^2(2x))$. For $\Re s>1$ the Mellin transform of $\Psi_f(x)-x/4$ on $[1,\infty)$, with measure $x^{-s-1}dx$, is
\[
 T(s)=-\frac1s\sum_\chi c_\chi\frac{L'}{L}(s,\chi)
                     -\frac1{4(s-1)}+U(s),
\]
where $c_\chi$ are the coefficients of \eqref{eq:identity}. The finitely many omitted Euler factors give $U$ holomorphic in $\Re s>0$. The principal pole cancels. Dirichlet GRH controls all four linear terms. At each zero $z$ of $L(s,\rho\eta)$ with $\Re z>1/2$, the residue of $T$ is $m_z/(3z)$, with $m_z$ its multiplicity.

Fix $1/2<\theta<1$. Equation~\eqref{eq:cubic} for a smaller exponent and the prime-power bound make $Cx^\theta-\Psi_f(x)+x/4$ nonnegative on $[1,\infty)$ for some $C$. Its Mellin transform is $C/(s-\theta)-T(s)$. If its abscissa of convergence were $>\theta$, Landau's theorem would force a singularity at that real point. The only possible pole there has strictly negative residue. A nonnegative integral approached from the right cannot have such a residue. The integral thus converges and is holomorphic on $\Re s>\theta$, excluding all poles there. Let $\theta\downarrow1/2$ and apply the functional equation. Conversely, the fixed-conductor explicit formula under RH for these finitely many Dirichlet and Hecke functions bounds their centered von Mangoldt sums by $O(\sqrt x\log^2(2x))$ (or $O_\epsilon(x^{1/2+\epsilon})$), giving \eqref{eq:cubic} by \eqref{eq:identity}; see also the Artin explicit formula in \cite{ng}.
\end{proof}
Equation~\eqref{eq:cubic} is the outstanding arithmetic estimate. The upstream covolume bound was built for a different prime weight. In particular the two events at the identity and at a 3-cycle of $S_3$ can have the same image in the abelianization and different values of $f$. The character identity exhibits exactly where a new estimate is required.

\section{Further research questions}
\begin{enumerate}
\item How small a finite set of primes can support a strict improvement on Rankin? What explicit increase in the natural density can be achieved by refining the common deletion? Positive density gain is established here; the optimal value remains open.
\item Is there a positive loss $\bar\alpha_{p^-}-\bar\alpha_p$ at every sufficiently large prime $p$, where $p^-$ is its predecessor? The proof at $11$ reduces strictness to finding a single interval on which the independent layers have a strictly larger optimum. A family of such intervals could quantify convergence from smooth to unrestricted ratios.
\item For the cyclotomic witnesses, which primes that divide the exact gcd support a nonzero odd component? The calculation rules out some primes universally. Establishing that a surviving prime occurs requires information about cocycles themselves.
\item How does the maximal sparse-place density decay when the prime cutoff grows with the source's parameter $s$? The fixed-cutoff proof tracks its dependence on the number of primes. Exact two-prime extremizers give specific candidates on which to test the remaining congruences.
\item Can the main lattice inequality be realized for the event in Theorem~\ref{th:cubic}, including its negative standard-twist coefficient and its vanishing untwisted coefficient? This would turn the single one-sided estimate into a zero theorem for an explicit non-Dirichlet function.
\end{enumerate}
\begin{thebibliography}{9}
\bibitem{grh} L. Sernau and collaborators. \emph{GRH for Dirichlet $L$-functions, in Lean 4}. \href{https://github.com/lsernau/grh-lean/tree/72fead102f663d5c7a407964a13dc0cf9d977207}{Pinned repository} and \upsrc{GUIDE.md}{reading guide}.
\bibitem{alpoge} L. Alp\"oge. \href{https://alpo.ge/remark.pdf}{\emph{Remarks on the proof}} (November 2026).
\bibitem{mcnew} N. McNew. \emph{On sets of integers which contain no three terms in geometric progression}. Mathematics of Computation \textbf{84} (2015), 2893--2910. \href{https://www.nathanmcnew.com/GPFsets.pdf}{Author's manuscript}, Sections 4--7.
\bibitem{riddell} J. Riddell. \emph{Sets of integers containing no $n$ terms in geometric progression}. Glasgow Math. J. \textbf{10} (1969), 137--146. \href{https://doi.org/10.1017/S0017089500000690}{DOI}.
\bibitem{packetgap} ProveIt. \prsrc{grh_transfers_density_gaps_and_zero_tests.zip}{\emph{Density gaps and zero tests}}. In particular, the \texttt{integers.tex} argument.
\bibitem{packetmc} ProveIt. \prsrc{mcnew_density_and_grh_transfers.zip}{\emph{A density separation for geometric progressions}}. See also its \texttt{frobenius.tex} on strictly negative binary tests.
\bibitem{packetcyc} ProveIt. \prsrc{grh_method_transfers_2026-11-26.zip}{\emph{Method transfers}, 26 November 2026}. See \texttt{cyclotomic.tex} and the question on exceptional coefficient primes in \texttt{report.tex}.
\bibitem{serrepoly} C. B. Khare, A. F. La Rosa and G. Wiese. \emph{Splitting fields of $X^n-X-1$ (particularly for $n=5$), prime decomposition and modular forms}. \href{https://arxiv.org/abs/2206.08116}{arXiv:2206.08116}.
\bibitem{ng} N. C. Ng. \emph{Limiting Distributions and Zeros of Artin $L$-Functions}. Ph.D. dissertation, University of British Columbia, 2000. \href{https://www.cs.uleth.ca/~nathanng/RESEARCH/phd.thesis.pdf}{Author's text}.
\end{thebibliography}
\end{document}
